\documentclass[12pt]{article}
\usepackage{graphicx} % Required for inserting images
\usepackage[margin=0.5in]{geometry}
\usepackage{amsmath, amssymb}
\usepackage{mathrsfs}


\usepackage{soul}


% Dr. David Brown Commands
\newcommand{\ds}{\displaystyle}
\newcommand{\R}{\mathbb{R}}
\newcommand{\N}{\mathbb{N}}
\newcommand{\Q}{\mathbb{Q}}
\newcommand{\C}{\mathbb{C}}


% Commands to get script r
\usepackage{calligra}
\DeclareMathAlphabet{\mathcalligra}{T1}{calligra}{m}{n}
\DeclareFontShape{T1}{calligra}{m}{n}{<->s*[2.2]callig15}{}
\newcommand{\scripty}[1]{\ensuremath{\mathcalligra{#1}}}

% Unit commands
\newcommand{\degree}{\text{°}}
\newcommand{\unitSpeed}{\frac{\text{m}}{\text{s}}}
\newcommand{\unitAcceleration}{\frac{\text{m}}{\text{s}^2}}

% Constant commands
\newcommand{\avogadro}{6.022\times10^{23}}

\newcommand{\boltzmannConstK}{1.381\times10^{-23}\frac{\text{J}}{\text{K}}}
\newcommand{\boltzmannConstInEV}{8.617\times10^{-5}\frac{\text{eV}}{\text{K}}}
\newcommand{\planckConst}{6.626\times10^{-34}\text{J}\cdot\text{s}}

\newcommand{\atmP}{1.01325\times10^5\,\text{Pa}}

% Known Value Commands
\newcommand{\electronMass}{9.109\times10^{-31}\,\text{kg}}
\newcommand{\protonMass}{1.673\times10^{-27}\,\text{kg}}

% Commands to easily write partial derivatives
\newcommand{\partDeriv}[2]{\frac{\partial #1}{\partial #2}}
\newcommand{\doublePartDeriv}[2]{\frac{\partial^2 #1}{\partial^2 #2}}


\title{Problem 7.42}
\author{Gabriel Decker}


\begin{document}


\maketitle
\begin{flushleft}


In this problem, we will consider the electromagnetic radiation inside a kiln at $T=1500\,\text{K}$.
We'll consider the kiln to have a volume of $V=1\,\text{m}^3$.
We will find the total energy of the radiation, plot the spectrum while varying photon energy $\epsilon$, and find the fraction of the energy in the visible spectrum (400 nm to 700 nm).\\~\\

The total energy can be determined by using the following equation from the book.
\begin{equation}
    \frac{U}{V}=\frac{8\pi^5(kT)^4}{15(hc)^3}
\end{equation}
We can also use the spectrum equation, as a function of photon energy $\epsilon$ from the book.
\begin{equation}
    u(\epsilon)=\frac{8\pi}{(hc)^3}\frac{\epsilon^3}{e^{\epsilon/kT}-1}
\end{equation}
We also found in problem 7.39 the total energy with a spectrum as a function of $\lambda$.
\begin{equation}
    \frac{U}{V}=8\pi hc\int_0^\infty\frac{1}{\lambda^5}\frac{1}{e^{hc/\lambda kT}-1}\,d\lambda
\end{equation}
If we want to get just the contribution to the energy from the visible spectrum, we can adjust the bounds of integration to the following.
\begin{equation}
    \frac{U_\text{vis}}{V}=8\pi hc\int_{400\,\text{nm}}^{700\,\text{nm}}\frac{1}{\lambda^5}\frac{1}{e^{hc/\lambda kT}-1}\,d\lambda
\end{equation}
We'll use all these things as we find the various desired quantities.\\~\\


\textbf{Part a:}\\
First, we'll find the total energy.
$$\frac{U}{V}=\frac{8\pi^5(kT)^4}{15(hc)^3}$$
$$\implies U=V\frac{8\pi^5(kT)^4}{15(hc)^3}$$
Now, we simply plug in our $V=1\,\text{m}^3$ and $T=1500\,\text{K}$.
$$\implies U=(1\,\text{m}^3)\frac{8\pi^5((\boltzmannConstK)(1500\,\text{K}))^4}{15((\planckConst)(2.998\times10^8\,\text{m/s}))^3}$$
$$\implies U=0.00383\,\text{J}$$\\~\\~\\



\textbf{Part b:}\\
The following is a plot of the spectrum as a functiuon of photon energy.\\
\includegraphics[width=0.6\linewidth]{energy_vs_spectrum.png}\\~\\~\\



\textbf{Part c:}\\
We will use the following equation to find the total energy from the visible spectrum.
We will divide this from the total energy to get the fraction from visible light.
$$\frac{U_\text{vis}}{V}=8\pi hc\int_{400\,\text{nm}}^{700\,\text{nm}}\frac{1}{\lambda^5}\frac{1}{e^{hc/\lambda kT}-1}\,d\lambda$$
This integral can be numerically evaluated using python.
$$\implies\frac{U_\text{vis}}{V}=2.136\times10^{-6}\,\text{J/m}^3$$
With $V=1\,\text{m}^3$, this implies the following.
$$\implies U_\text{vis}=2.136\times10^{-6}\,\text{J}$$
The fraction of energy density from visible light is the following.
$$\frac{U_\text{vis}/V}{U/V}=\frac{U_\text{vis}}{U}$$
$$\implies\frac{U_\text{vis}/V}{U/V}=\frac{2.136\times10^{-6}\,\text{J}}{0.00383\,\text{J}}$$
$$\implies\frac{U_\text{vis}/V}{U/V}=0.000558\%$$





\end{flushleft}

\end{document}
